\begin{thebibliography}{9}
\bibitem[Source(2026)]{source}
\emph{The $(3,4,\infty)$ modular family of 2-tori, completed at its three
special points, is a complex structure on $S^6$}. Manuscript circulated
by Levent Alp\"oge; Engel attributes the manuscript to Claude.
Public 108-page version inspected 12 September 2026.
Setup, \S2 Lemmas 2.4--2.7, \S3 period conditions, \S4 cusp,
Proposition 5.6, \S7 topology. \url{https://alpo.ge/s6.pdf}.
The global sphere claim is not independently certified here.
\bibitem[Engel(2026)]{engel}
Engel, P. (2026). \emph{Complex structures on $S^6$}.
Twenty-five-page exposition, introduction and initial constructions.
\url{https://philip-engel.github.io/S6.pdf}.
\bibitem[S6 workbench(2026)]{s6repo}
\emph{Yang--Mills, S6 and Navier--Stokes research}. S6 key advances and
attempts, frozen 6 September 2026. Read at revision
\texttt{8cfbd18f3e398c533e1acabfba98b4a29975000a}.
Key-advances \S\S1--3 distinguish inverse periods, deformations and attachments.
\url{https://github.com/KokunoYumeto/yang-mills-interacting-workbench}.
\bibitem[Sijsling and Voight(2014)]{sv}
Sijsling, J., \& Voight, J. (2014). On computing Bely\u\i\ maps.
\emph{Publications math\'ematiques de Besan\c con}, 2014(1), 73--131.
Introduction and \S1: permutation triples, covers, dessins and genus;
see also the authors' 2023 errata for the precise cyclic-order definition.
\url{https://arxiv.org/abs/1311.2529}.
\bibitem[Elsholtz and Tao(2013)]{et}
Elsholtz, C., \& Tao, T. (2013). Counting the number of solutions to the
Erd\H{o}s--Straus equation on unit fractions.
\emph{Journal of the Australian Mathematical Society, 94}(1), 50--105.
Section 2, Propositions 2.2 and 2.6.
\url{https://arxiv.org/abs/1107.1010}.
\end{thebibliography}
